\documentclass[11pt]{article}
\usepackage{amsmath,amssymb,amsthm}
\newtheorem{theorem}{Theorem}
\newtheorem{corollary}{Corollary}
\title{Independent exterior signs under a fixed polynomial sum of squares}
\author{}\date{}
\begin{document}\maketitle

\paragraph{Status.}
This is a complete written single-layer construction, with exact finite
rational checks. It is not Lean-checked. The four-child corollary is a
compatible two-level polynomial construction, but it does not prescribe
independent signs at its two child splits. No recursive compiler or
asymptotic head bound is claimed.

Fix distinct real query centers $\tau_1,\ldots,\tau_D$, with $D\geq1$.
For a sign word $\sigma\in\{-1,1\}^D$ and a positive parameter $\eta$,
put
\[
 Z_\sigma(t)=\prod_{j=1}^D(t-\tau_j+i\sigma_j\eta),\qquad
 C_\sigma=(1+i)Z_\sigma=A_\sigma+iB_\sigma,
 \qquad R_\sigma=A_\sigma^2-B_\sigma^2.
                                                               \tag{1}
\]
The real polynomials $A_\sigma,B_\sigma$ are both monic of degree $D$.
In particular $\deg R_\sigma\leq2D-1$. Their sum of squares is fixed:
\[
 A_\sigma^2+B_\sigma^2
       =2\prod_{j=1}^D\bigl((t-\tau_j)^2+\eta^2\bigr)
       =:S_\eta(t),                                           \tag{2}
\]
independently of $\sigma$. Thus the binary choices in (1) change the
real difference of the squares while preserving their sum.

\begin{theorem}[Fixed-norm spectral signs]
Let $P$ be any fixed real polynomial with $P(\tau_j)>0$ for every
$j$. For all sufficiently small $\eta>0$, uniformly over all sign
words $\sigma$, the polynomials in (1) satisfy
\[
 -\sigma_j R_\sigma(\tau_j+\eta)>0,
 \qquad
 -\sigma_j(R_\sigma/P)'(\tau_j)>0
                         \quad(1\leq j\leq D).                \tag{3}
\]
\end{theorem}
\begin{proof}
Fix $k$ and write
\[
 G_k(t)=\prod_{j\ne k}(t-\tau_j+i\sigma_j\eta),\qquad
 g_k=\prod_{j\ne k}(\tau_k-\tau_j)\ne0.
\]
Since $C_\sigma^2=2iZ_\sigma^2$, one has
$R_\sigma=-2\operatorname{Im}(Z_\sigma^2)$. At $t=\tau_k+\eta$,
\[
 R_\sigma(t)=-4\sigma_k\eta^2\operatorname{Re}(G_k(t)^2).
                                                               \tag{4}
\]
The last real part tends to $g_k^2>0$ as $\eta\to0$.
This proves the first inequality.

At the center $t=\tau_k$, direct differentiation gives
\[
 \begin{aligned}
 R_\sigma(\tau_k)&=2\eta^2\operatorname{Im}(G_k(\tau_k)^2),\\
 R_\sigma'(\tau_k)
 &=-4\sigma_k\eta\operatorname{Re}(G_k(\tau_k)^2)
       +4\eta^2\operatorname{Im}(G_k(\tau_k)G_k'(\tau_k)).
 \end{aligned}                                                 \tag{5}
\]
The coefficients of $G_k$ are polynomial in $\eta$ with real
constant terms. Their real parts are even in $\eta$ and their
imaginary parts are odd. Consequently (5) implies
\[
 R_\sigma(\tau_k)=O(\eta^3),\qquad
 R_\sigma'(\tau_k)=-4\sigma_k g_k^2\eta+O(\eta^3).
\]
The quotient rule now yields
\[
 (R_\sigma/P)'(\tau_k)
       =-\frac{4\sigma_k g_k^2}{P(\tau_k)}\eta+O(\eta^3),
                                                               \tag{6}
\]
which has the required strict sign. There are only finitely many
pairs $(\sigma,k)$, so a common positive upper bound on $\eta$
works for all of them. This uniformity is for each fixed $D$, query
set, and $P$; no dimension-independent numerical bound is asserted.
\end{proof}

\begin{corollary}[Four compact real-rooted child blocks]
Let $U$ be monic of degree $M>D$, with $M$ distinct real roots in
an interval $(a,b)$ lying strictly below every $\tau_j$. There is an
$\eta>0$ such that the following holds. For every sufficiently small
$\delta>0$, uniformly over all $\sigma$, the four polynomials
\[
 U+\delta A_\sigma,\quad U-\delta A_\sigma,\quad
 U+\delta B_\sigma,\quad U-\delta B_\sigma                  \tag{7}
\]
are monic and have $M$ distinct real roots in $(a,b)$. Let
\[
 G_A=U^2-\delta^2 A_\sigma^2,\quad
 G_B=U^2-\delta^2 B_\sigma^2,\quad
 P_\delta=U^2-\delta^2S_\eta/2.
\]
Their mean $P_\delta=(G_A+G_B)/2$ is independent of $\sigma$, both
pairs in (7) have common mean $U$, and the root transforms satisfy
\[
 \sigma_j\left(\frac{G_A'(\tau_j)}{G_A(\tau_j)}
                   -\frac{G_B'(\tau_j)}{G_B(\tau_j)}\right)>0
                      \quad(1\leq j\leq D).                  \tag{8}
\]
Each logarithmic derivative is the Stieltjes transform of the roots
of its polynomial, with all weights equal to one and multiplicities
counted.
\end{corollary}
\begin{proof}
Apply the theorem to $P=U^2$ and fix a sufficiently small $\eta$.
Choose $M$ disjoint closed intervals around the roots of $U$,
contained in $(a,b)$, with nonroot endpoints. For small $\delta$,
every polynomial in (7) has the sign of $U$ at each endpoint.
The intermediate value theorem gives a root in every interval;
monicity and degree $M$ imply that these are all the roots and are
simple. Finite uniformity over $\sigma$ gives a common $\delta$
bound. Different polynomials in (7) may share roots.

Direct expansion gives
\[
 \frac{G_A'}{G_A}-\frac{G_B'}{G_B}
     =-\delta^2\frac{R_\sigma'P_\delta-R_\sigma P_\delta'}
                         {G_A G_B}.                            \tag{9}
\]
At each query the denominator is positive. As $\delta\to0$,
the numerator in (9) tends to
$R_\sigma'(U^2)-R_\sigma(U^2)'$, which by (3) has sign
$-\sigma_j$. Thus (8) follows for sufficiently small positive
$\delta$, uniformly over the finite sign family. The logarithmic
 derivative identity follows by differentiating the product of
linear factors, including repeated factors if present.
\end{proof}

\paragraph{The remaining recursive gap.}
The corollary prescribes $D$ root bits using $4M$ unit-weight poles.
It does not give a ratio of queries to poles growing with $M$.
At the centers, as $\eta\to0$, both $A_\sigma'$ and $B_\sigma'$
have the same leading value $g_k$. The two child split transforms
in (7) therefore have leading signs fixed by $g_k$ and $U$, rather
than independently prescribed child words. An induction would need
additional freedom that simultaneously preserves the common sums
of squares and realizes the required descendant differences.

Factoring an arbitrary child remainder as a difference of squares
does not by itself close this gap: the corresponding sum of squares
changes its ancestor's mean. Nor does counting the $2^D$ spectral
factors prove coverage of more than these $D$ explicitly read bits.
The constant-remainder collapse in
\texttt{underloaded-equal-weight-split.tex} is avoided here because
$R_\sigma$ is allowed degree $2D-1$; that observation is only local.

\paragraph{Exact finite verification.}
\texttt{spectral-factor-root-signs-check.py} takes rational centers,
width, poles, and perturbation size. For every $1\leq D\leq8$ and
every sign word, it checks the norm identity (2), the degree drop,
both strict signs in (3), isolating endpoint certificates for the
four polynomials (7), their products and common means, and the
identity and strict signs in (9). There are $510$ complete rational
certificates. No numerical root solver or floating-point sign test
is used, and no recursive sign-coverage assertion is tested.
\end{document}
